\section{DISCUSSION}
\label{sec:discussion}

The results support the two-part hypothesis, but only as stage-level evidence.
They identify where this implementation is likely constrained and what kind of
future end-to-end benchmark is needed; they do not yet prove superiority over
other PDF translation systems.

\subsection{H1: What Kind of Error Bounds the Pipeline}
\label{subsec:disc-structure}

The parsing results make a useful distinction between benign structure changes
and unrecoverable loss. Over-segmentation is visible in the gap between strict
COCO 1-1 localisation and the granularity-invariant $N$--$M$ matcher, but it is
mostly compatible with the downstream contract: fragments can be gathered,
ordered and translated as long as their text and geometry remain present. A
box-exact metric alone would therefore select against a design choice that is
helpful for translation, especially around sparse formula blocks.

The errors that matter are different. A missed region supplies no text to the
translator; a missed formula or table removes structured content; and OCR errors
turn into fluent but wrong source strings. These failures are concentrated in
specific slices---degraded historical pages, CJK OCR, three-column formulas and
table coverage---which makes them actionable. The appropriate response is not a
uniformly larger translator, but targeted parser/OCR improvements and a review
checkpoint exactly where loss first occurs.

The translation experiment gives the other half of H1. The five evaluated LLMs
do differ, but their mean COMET-DA scores occupy a narrow band compared with the
variation across target languages. That means model choice is still worth doing,
but it is a smaller controllable lever than the document-specific stages once a
competitive hosted model is used. The strongest next test of H1 is therefore
direct: improve the failing parser slices and compare the end-to-end effect
against switching translation models.

\subsection{H2: Layout Belongs to Reconstruction}
\label{subsec:disc-constraint}

The length-compliance results show why layout preservation should not be
delegated to the translator. A character budget is weakly followed and mostly
tracks target-script density: CJK output is too short by character count even
when adequacy is high, while some Latin-script targets violate by expansion.
Because the correlation with COMET-DA is effectively zero, compliance is not a
useful proxy for semantic quality.

More importantly, a character count is not the quantity the PDF cares about. The
actual layout question is whether rendered text, in the chosen font, at a local
position, collides with neighbouring elements or remains readable after
shrinking. Only the reconstruction stage has the page geometry and pixels needed
to answer that. The pipeline's size and colour decisions therefore belong in
reconstruction: size is changed only when rendered overflow would collide, and
appearance is sampled from the source page instead of assumed to be black text on
white background.

\subsection{Metric Design}
\label{subsec:disc-metrics}

The parser evaluation intentionally uses several metrics because each one
answers a different failure question. Localisation asks whether the region is
available at all; classification asks whether the right policy is applied; OCR
asks whether the translator receives the right string; reading order asks
whether context is preserved; formula CDM separates rendered mathematical
equivalence from LaTeX spelling; and table-content scoring acknowledges that
cell structure cannot be judged without cell-level ground truth. Reported
together, these metrics describe how an error would propagate through the
pipeline rather than compressing incompatible failures into one score.

\subsection{Limitations and Next Steps}
\label{subsec:disc-limitations}

Several limits remain. First, the study is not an end-to-end decomposition of
final PDF quality. It measures the parser and translator separately, then uses
the pipeline mechanism to argue which errors are recoverable. Second, final
layout fidelity is shown qualitatively but not scored with an output metric such
as overflow rate, font-shrink rate, colour error, visual similarity or
page-count preservation. Third, no direct baseline is run against systems such
as PDFMathTranslate, docutranslate or RetainPDF under one shared input and output
contract. Fourth, the cell stage is now scored, but only for geometry and only on cropped
scientific tables: no public corpus provides full pages, cell boxes that cover blank
cells, and annotations free of the over-segmentation that affects earlier datasets. Cell
geometry and table detection are therefore measured on different corpora, and logical
structure is not measured at all, because the parser does not produce it. Finally,
the experiments focus on \texttt{en}$\rightarrow$\texttt{xx}; other source
languages and document genres remain open.

The next benchmark should therefore be an annotated set of real PDFs containing
flowing text, tables, formulas, scanned pages and born-digital pages. It should
score both adequacy and final artifact quality, compare direct PDF-translation
systems and conversion-based pipelines under the same reconstruction contract,
and test whether improving parser coverage changes the final PDF more than
changing the translation model. That is the measurement needed to turn the
stage-level ordering argued here into an end-to-end claim.
